% AC-G138-quartic-model. Complete author manuscript; independent review pending.
% Second residue only. First quartic label Psi(2)=i; G137 is checked.
\section{The second quartic residue on its actual host}
\label{sec:cubic-quartic-second-residue-actual-model}

This section constructs the actual host deformation rings and modular
ideals for the second weight residue of the first quartic character.
The ordinary branch is newly chosen through a weight-\(267\) modular
form, and the native ideals are constructed on its actual residual
tuple.  Its relation to the first residue is
\[
 \bar\gamma_{133}
      =\bar\gamma_{267}\tau_p,\qquad
                    \tau_p=\bar c_p^{134}.
\]
The finite smooth \(p\)-inertia characters remain the cubic triple,
but the algebraic finite unit-torus character is recomputed below.
No modular ideal from AC-G136 is identified with this new one.
The auxiliary rational form is used after its extra finite character
is killed; it is never declared to be an original rational
deformation point.

\paragraph{Original domain and retained inputs.}
Put
\[
 \begin{gathered}
 p=269,\quad K=\mathbb Q(\sqrt{-3}),\quad
 \eta=\chi_{-3},\quad N=10,\\
 \psi_5(2)=i,\quad i\equiv82\pmod p,\quad
 e\equiv133\pmod{268},\quad e>0.
 \end{gathered}
 \tag{G138.1}
\]
Here \(\psi\) is the character modulo ten induced by the
primitive quartic character \(\psi_5\) of conductor five.
Write \(\Psi=\operatorname{DirG}(\psi_5)\), using arithmetic
Frobenius for this finite avatar and for the actual cyclotomic
character \(c_p\); thus
\(c_p(\operatorname{Fr}^{\rm geom}_w)=q_w^{-1}\).
For \(j\in\{1,2\}\), retain the original coefficient DVR
\(\mathcal O_0\), the actual cubic selected-component ring
\(S_j\), and its model, invariant subring and cotangent length.
The exact original \(\mathrm{GD}/\mathrm{Dat}\), local,
coefficient, and marked-element hypotheses remain those of
AC-G65--AC-G68.  In particular the old modularity supply used
in AC-G128 retains \(e_{\mathcal O_0}=1\).

Retain the published modular attachment, ordinary control,
endoscopic transfer, local--global compatibility, and cyclic
base-change inputs of AC-G78--AC-G80 and AC-G127; the arithmetic
class-field and duality inputs of AC-G74 and AC-G90--AC-G100;
the actual native \(\mathrm{LevEx}\),
\(\mathrm{ProModularRingExists}\), \(\mathrm{LamEx}\), and
\(\mathrm{ModU}\) supplies; and the original generic algebraic
\(\mathrm{CohInp}\), \(\mathrm{LemmaP}\),
\(\mathrm{LemmaM}\), and \(\mathrm{LemmaM'}\) inputs.
The exact original-field descent and counting inputs are those
of AC-G110--AC-G113.  The additional finite arithmetic input
is precisely AC-G137-ramified-quartic-classes for
\(d=-2959,538,-5918\), together with the already checked AC-G133
and AC-G134 base sector.  All their stated source scopes remain.

\paragraph{Host notation and original model.}
The auxiliary prime is \(q_{\rm aux}=16141\).  Put
\[
 F_1=K\mathbb Q(\sqrt{-11}),\qquad
 E=\mathbb Q(\sqrt2),\qquad
 H_{\rm ab}=F_1E.
 \tag{G138.2}
\]
Write \(E_0=\operatorname{Frac}(\mathcal O_0)\), and denote the
original model by
\[
 (A_0,\mathcal A,x,q_{\rm ring}:A_0\longrightarrow S_j,\mathcal Q),
 \qquad I_0=\mathcal Q.\mathrm{ideal}.
 \tag{G138.3}
\]
Let \(\mathscr B_0\) be its complete degree-zero ring and \(L_j\)
its weight-minus-one cotangent length.  Put
\(\gamma_e=\Psi c_p^e\) and
\(\bar\gamma=\bar\Psi\bar c_p^{133}\).
The original \(\mathcal O_0\), all these rings, the selected
local prime, and their length normalization are fixed throughout.
In particular the selected local prime is the one belonging to
this original \(e\), not a component transported from AC-G136.

\begin{theorem}[AC-G138-quartic-model: actual second-residue quartic host and original model]
\label{AC-G138-quartic-model}
Under the retained inputs above, including the exact class
conclusion of AC-G137, there is a finite CM host \(H/F\),
\(F=H^+\), with its full local Galois groups, a finite host
coefficient DVR \(\mathcal O\), and actual tuple rings
\[
 B\subset A\subset R,\qquad
 \theta=(\bar\chi^j\bar\gamma,1,
                         \bar\chi^j\bar\gamma^{-1}).
 \tag{G138.4}
\]
For fixed \(\mathcal O_0,j\), the host, these rings, and the
native modular ideals may be chosen once for every original
weight \(e\) in (G138.1).  These are new second-residue quartic
data.  The initial fields \(F_1,E,T\) may be kept as in AC-G136,
but the ordinary branch, the later prime \(q\), the final host,
and the native ideals are obtained by the construction below.

For the actual native map \(\Phi_3:\mathcal B_3\to B\) and
its proper modular ideal \(J_3\), set
\(I=\Phi_3(J_3)B\).  There is an actual uniform even ideal
\(M_E\subset B\) such that, for \(I_{\rm off}=I\cap M_E\),
\[
 \begin{aligned}
 I&\subset\mathfrak q\cap B
       &&\text{for every Type I prime of }A,\\
 I_{\rm off}&\subset\mathfrak q\cap B
       &&\text{for every off-triangular prime of }A,\\
 \sqrt{I_{\rm off}+\lambda_\theta}&=\mathfrak m_B.
 \end{aligned}
 \tag{G138.5}
\]
Here \(\lambda_\theta\) is the actual joint-residual-inertia
weight ideal, including the host uniformizer.
The first inclusion is the precise point-level condition
(G110.2).  In particular
\(\mathscr B_0\) is finite over \(\mathcal O_0\), every original
two-sided characteristic-zero point is Type I and rigid, and
the original model (G138.3) satisfies all five literal
conditions \(\mathrm{M1}\)--\(\mathrm{M5}\).

For
\[
 \begin{split}
 \alpha_j&=\chi^{-j}\Psi^{-1}c_p^{-e}|_{G_K},\qquad
 \mathcal D_j=(E_0/\mathcal O_0)(\alpha_j),\\
 \mathcal S_j&=\ker\bigl(H^1(G_K,\mathcal D_j)\\
 &\hspace{3em}\longrightarrow
 \prod_{v\mid p}H^1(G_{K_v},\mathcal D_j)
 \times\prod_{v\nmid p}H^1(I_v,\mathcal D_j)\bigr),\\
 s_j&=\operatorname{length}_{\mathcal O_0}(\mathcal S_j)/
                                  e_{\mathcal O_0},
 \qquad Y=10(e+1),\qquad b=121^e+1+121^{-e},
 \end{split}
 \tag{G138.6}
\]
these lengths are finite.  At geometric Frobenius \(\sigma\)
of the degree-two place of \(K\) over \(11\), there is a finite
family \(X_j\) of actual Type I representations, pairwise
nonisomorphic on \(G_K\), with
\[
 s_j\le L_j/e_{\mathcal O_0}
 \le2\sum_{\rho\in X_j}
            v_{269}(\operatorname{tr}\rho(\sigma)-b),
 \qquad \#X_j\le C_1Y^{A_1}.
 \tag{G138.7}
\]
The constants are precisely those of AC-G112-count.
For every nonempty \(J_0\subset\{1,2\}\), one nonzero integer
\(\mathcal N_{e,J_0}\) satisfies
\begin{align}
 \sum_{j\in J_0}L_j/e_{\mathcal O_0}
      &\le2v_{269}(\mathcal N_{e,J_0}),\notag\\
 \log|\mathcal N_{e,J_0}|
      &\le5C_1\#J_0\,Y^{A_1+1}\log11,\notag\\
 \left(\sum_{j\in J_0}L_j/e_{\mathcal O_0}\right)\log269
      &\le10C_1\#J_0\,Y^{A_1+1}\log11.
 \tag{G138.8}
\end{align}
The same estimates hold with \(s_j\) in place of the normalized
cotangent lengths.  The two cubic labels may use separate hosts;
the integer uses their common original-field classifier.
No host coefficient ramification index rescales these lengths.
\end{theorem}

\begin{proof}
\proofstep{1. The auxiliary cusp form at the new prime.}
The exact integer certificate accompanying this proof checks
primality and
\[
 \begin{gathered}
 q_{\rm aux}-1=60p,\qquad q_{\rm aux}\equiv1\pmod{15},\qquad
 6\text{ primitive modulo }q_{\rm aux},\\
 6^{1084}\equiv269\pmod{q_{\rm aux}},\qquad1084\equiv8\pmod p.
 \end{gathered}
 \tag{G138.9}
\]
Choose \(\theta_0(6)=\zeta_p\), let \(\Theta_0\) be its
arithmetic Galois avatar, and let \(T/\mathbb Q\) be its real
cyclic field of degree \(p\).  Put \(\xi_0=\eta\psi_5^{-2}
=\chi_{-15}\).  Then \(\xi_0(q_{\rm aux})=1\), and
\(\theta_0(p)=\zeta_p^8\ne1\).

We replay the prime-dependent steps of AC-G127, rather than
apply its fixed-3229 statement at a different prime.
The primitive odd character \(\xi_0\theta_0\) has conductor
\(15q_{\rm aux}\).  Katz's integral Eisenstein moment gives
\[
 \frac{L(1-k,\xi_0\theta_0)}2
 \equiv(1-\xi_0(q_{\rm aux})q_{\rm aux}^{k-1})
             \frac{L(1-k,\xi_0)}2\equiv0\pmod\varpi
 \tag{G138.10}
\]
for every odd \(k>2\).  The smoothing factor reduces to
\(1-2^k\), a unit because \(2^{134}=-1\pmod{269}\).
This also proves the required integrality without dividing by
\(k\).  The primitive cusp-constant formula makes every cusp
constant a unit multiple of this value or zero.  The reduction
is a nonzero cusp form since its coefficient of \(Q\) is one.

On the fine compact curve for
\(\Gamma_1(15q_{\rm aux})\cap\Gamma(4)\), its diamond
subgroup \(P\) of order \(p\) acts freely on every geometric
fiber, including cusps.  Indeed the squarefree-cusp stabilizer
calculation of AC-G127 gives only the trivial diamond, up to
sign, and full level four removes elliptic-curve automorphisms.
The quotient is a finite etale torsor on the compact curve.
Descending the character line makes its degree zero and its
reduction trivial.  Serre duality and
\(\omega^2=\Omega^1(C)\) give
\(H^1(\omega^k\mathcal L_{\theta_0}(-C))=0\) on the
special fiber for \(k>2\), and hence the actual integral cusp
lift.  The complementary diamond group now has order \(480\),
and the final full-level-four group has order \(48\).
Both are units at \(p\).  These are the only averaging
denominators; no order-\(p\) character projector is used.
The full Hecke eigenlift includes the three bad roots and the
ordinary root at \(p\), exactly as in AC-G127.

Take \(e_*=133\), \(k_*=267\).  The actual primitive twist is
\[
 \begin{split}
 g&\in S_{267}(15q_{\rm aux},\xi_0\theta_0),\\
 f_*&=g\otimes(\eta\psi_5^{-1}\theta_0^{-1/2}),\\
 V&=R_{f_*}c_p^{133}
    =R_g\eta\Psi\Theta_0^{1/2}c_p^{133}.
 \end{split}
 \tag{G138.11}
\]
The source geometric-Frobenius convention gives
\(\det R_g=\eta\Psi^2\Theta_0^{-1}c_p^{-266}\).
Thus \(\det V=\eta\) and
\[
 \bar V^{\rm ss}=\eta\bar\gamma\oplus\bar\gamma^{-1}.
 \tag{G138.12}
\]
The factor \(\eta\) in (G138.11) is retained before restriction
to \(K\).  The primitive level of \(f_*\) is exactly
\(75q_{\rm aux}^2\), its nebentype is \(\eta\), and it is
crystalline ordinary at \(p\).  Its exact away inertia is
\(1\oplus\eta\) at three, \(\Psi\oplus\Psi^{-1}\) at five,
\(\Theta_0^{1/2}\oplus\Theta_0^{-1/2}\) at \(q_{\rm aux}\),
and trivial elsewhere; monodromy is zero.  These conductor
assertions use primitive local factors, not an imprimitive
coefficientwise Dirichlet twist.  In particular \(V\) and its
entire fixed-tame-level ordinary determinant are unramified at two.
This is an actual auxiliary seed, not an original
\(\mathrm{Dat.IsPt}\).

\proofstep{2. The initial host and the whole ordinary branch.}
Write \(K'=\mathbb Q(\sqrt{-11})\) and \(F_*=F_1T\).
The coprime discriminants \(-3,-11\) give
\(\operatorname{Disc}(F_1)=33^2\) and
\(\operatorname{Disc}(F_1^+)=33\).  Hence \(F_1/F_1^+\)
is unramified at every finite place.
The exact symbols are
\[
 (-11/3)=(-11/5)=(-11/p)=(-11/q_{\rm aux})=1,
 \qquad (-3/p)=-1.
 \tag{G138.13}
\]
Thus all ramified places of the block split relatively, and
\(p\) splits in the source imaginary quadratic field \(K'\).
Both \(K\) and \(K'\) are unramified quadratic locally at two;
relative splitting there is not required for an unramified block.
Since \(\theta_0(p)\ne1\), every completion of \(F_*\) over
\(p\) is the unramified extension of \(\mathbb Q_p\) of
degree \(2p=538\).  The CM extension \(F_*/F_*^+\) is
finite-unramified and relatively split over \(p,3,5\).

The restriction of \(V\) to \(K\) is irreducible by the
ordinary inert-prime argument of AC-G127.  The next quadratic
restriction cannot split it: that would give a self-twist
ramified at eleven, where \(V\) is unramified.  The degree-\(p\)
restriction to \(F_*\) stays irreducible, because a character
self-twist of a rank-two representation has square one.

Choose an actual finite torsionfree domain ordinary branch
through this newly obtained \(f_*\) of weight \(267\), over the
half-weight variable \(W\), with
dense positive integer weights \(w\equiv133\pmod{268}\).
The normalized trace has determinant \(\eta\), and
\(1+W=(1+p)^w\) at such points.
The residual \(\bar\gamma\) is now
\(\bar\Psi\bar c_p^{133}\); this is the weight component used in
the actual branch and its later ordinary module.
The tame nebentype is fixed even though the diamond order is
divisible by \(p\): finite diamond values in the domain are
algebraic constants, and their embedding at the characteristic-zero
point is injective.  Thus the domain map factors through the
specified character.  No integral projector dividing by \(p\)
is used.

Saha's actual specialization theorem, with the retained hypotheses
of AC-G120, fixes the finite inertia and \(N=0\) on a dense
arithmetic component.  In particular the fixed
\(\Theta_0^{\pm1/2}\) inertia becomes trivial after restriction
to \(F_*\).  At three impose the unramified host type;
at five retain the exact distinct old endpoint characters
\(\Psi,\Psi^{-1}\).  At \(q_{\rm aux}\) impose maximal compact
level directly on the host, after the character has been killed.
At two the branch is already unramified from its Galois quotient.
The cubic character is unramified away from \(p\), including
at three: \(-269\in1+9\mathbb Z_3\), so
\(\exp(\log(-269)/3)\) is a cube root of \(-269\) in
\(\mathbb Q_3\).  Away from \(3p\), the unit Kummer
polynomial is etale.  Thus twisting by \(\chi^j\) adds
no further old away-inertia type.
With source imaginary field \(K'\), parallel archimedean pairs
are repeated \(2p\) times.  This even multiplicity, the relative
splitting at the ramified places, and the unramified nonsplit
clause give the literal Thorne 4.5 hypotheses.
The packet is ordinary by its actual local comparison.

Here is the exact unit-torus summand.  At every coefficient-prime
place let \(T_v(\mathcal O_v)\) be the full diagonal unit torus,
with pro-\(p\) subgroup \(T_v^1\), and put
\(\mu_*=(132,0,-132)\).
Define the finite character \(\delta_*\) on
\(T_v(\mathcal O_v)/T_v^1\) by restricting
\[
             (\chi^j,1,\chi^j)\,w_0(\mu_*)^{-1}
\]
to Teichmuller representatives.  Lowest-weight evaluation and
the canonical Jacquet map are equivariant for the \emph{full}
unit torus: at weight
\(\mu(w)=(w-1,0,1-w)\), the scalar ordinary unit character is
the smooth ordered unit character times \(w_0(\mu(w))^{-1}\).
Thus \(\delta_*\) is its actual finite part, not a projector
constructed from the pro-\(p\) action alone.  For
\(w\equiv133\pmod{268}\), the parallel algebraic weight difference
on residue units factors through their norm to \(\mathbb F_p^\times\);
its exponent is divisible by \(268\).  The same \(\delta_*\)
therefore serves the whole dense progression.

For completeness, let \(f=538\), let the increasing ordinary
slopes be \(s=(-fw,0,fw)\), and use the positive torus element
\(\operatorname{diag}(\varpi_v^2,\varpi_v,1)\).
The explicit normalized operator in AC-G79 has valuation
\[
                  2f-2fw+2f(w-1)=0 .
\]
For a different ordering \(\sigma\), its excess valuation is
\(\sum_{i=1}^3(3-i)(s_{\sigma(i)}-s_i)>0\), by strict
rearrangement.  Hence the ordinary line is exactly the ordered
line used to define \(\delta_*\).
From the start impose the commuting integral projectors
\(e_{\rm aux}e_{\rm old}e_{\delta_*}\): the prime-to-\(p\)
neat projector, the old compact types, and this finite
full-unit-torus character.
The old type at five distinguishes \(\Psi\) and \(\Psi^{-1}\);
it is not a common scalar type.
All commute with trace transitions, the ordinary operator,
and the later finite group at \(q\).
The free inverse-trace module, exact weighted-norm specialization,
and arbitrary-base-change Cayley--Hamilton argument of
AC-G78--AC-G80 therefore give the actual ordinary support,
not merely an interpolated characteristic polynomial.
These are applications with the preceding host and local
hypotheses rechecked.

\proofstep{3. A real quadratic field inert at the coefficient prime.}
The field \(E=\mathbb Q(\sqrt2)\) is ramified only at two,
where \(F_*\) and the whole ordinary branch are unramified.
Thus \(E\not\subset F_*\).
Its units modulo sign are generated by \(1+\sqrt2\), and
its class number is one.  The latter follows from the real
quadratic Minkowski bound \(\sqrt8/2<2\).
For the unit assertion, a positive unit between one and
\(1+\sqrt2\) would have its integral \(\sqrt2\)-coefficient
of absolute value at most one; its norm equation gives no
such unit.  Put
\[
 u=(1+\sqrt2)^2=3+2\sqrt2,
 \qquad M=E(\mu_p,u^{1/p}),\quad
 P=\operatorname{Gal}(M/E(\mu_p)).
 \tag{G138.14}
\]
The norm-one unit \(u\) generates a subgroup of index prime
to \(p\) in the full unit group.
It is not a \(p\)th power in \(E\), nor in \(E(\mu_p)\):
norming a putative root and using the degree prime to \(p\)
would give a root in \(E\).
Thus \(P\simeq C_p\); the argument below would also be
automatic if it were trivial.
The field \(M\) is Galois over \(\mathbb Q\), ramified only
at \(2,p\), with action \(\eta_E\bar c_p\) on \(P\).
Every nontrivial quadratic subfield of \(F_1\) ramifies at
three or eleven, so \(F_1\cap M=\mathbb Q\).
Also \(T\cap F_1M=\mathbb Q\), since \(T\) ramifies at
\(q_{\rm aux}\) and \(F_1M\) does not.  Therefore \(G_{F_*}\)
surjects onto \(\operatorname{Gal}(M/\mathbb Q)\), in particular
onto \(P\).

Choose \(h\in I_2\subset G_{F_*}\) with \(\eta_E(h)=-1\).
For the actual normalized branch trace \(\mathcal T\),
\[
 \mathcal T(ha)=\mathcal T(a)\quad(a\in G_{F_*}),\qquad
 c_p(h)=\chi(h)=\Psi(h)=1.
 \tag{G138.15}
\]
Choose \(t\in I_p\subset G_{F_*}\) with
\(c_p(t)=-(1+p)\) and \(\chi(t)=1\).
The coprime tame orders \(268,3\), together with the full wild
cyclotomic image of an unramified local host, give this choice.
Although \(E\) is inert at \(p\), it is unramified there;
hence \(\eta_E(t)=1\), as is \(\Psi(t)\).
Since all branch exponents are odd,
\[
             \mathcal T(t)=-((1+W)+(1+W)^{-1}).
 \tag{G138.16}
\]
Both \(h,t\) invert \(P\); their product centralizes \(P\)
and its square lies in \(P\).
For a lift in \(G_{F_*}\) of \(a\in P\), conjugating \(h\)
changes this square by \(a^4\).  Fourth powering is invertible
on \(P\), so choose the lift to obtain an actual \(g=h't\)
with \(g^2|_M=1\).  Cyclicity of trace preserves (G138.15)
for the conjugated \(h'\), hence preserves (G138.16).

Use this single actual element in AC-G78's exact-root argument,
or its finite-normalization form AC-G84--AC-G85, adjoining
\(F_*,M\), the residual fields, and \(c_p\bmod p^2\) to the
one finite quotient.  It gives a good rational prime \(q\),
completely split in \(F_*\), and an actual ordinary support
point with the required full root chain, such that
\[
 \begin{gathered}
 q\equiv-1\pmod p,\quad v_p(q+1)=1,\quad
 \chi(\operatorname{Fr}_q)=\Psi(\operatorname{Fr}_q)=1,\\
 \eta_E(q)=-1,\qquad
 u\text{ is a }p\text{th power in }\mathbb F_{q^2}^{\times}.
 \end{gathered}
 \tag{G138.17}
\]
The geometric convention gives
\(q^{-1}\equiv-(1+p)\pmod{p^2}\), which proves the exact
valuation in (G138.17).
The square-one prescription makes the inert \(E\)-prime split
in the Kummer extension, proving the unit assertion.
All conditions came from the one element \(g\); no independent
Chebotarev prescriptions were assumed compatible.

\proofstep{4. Actual finite type, the ray cover, and the local field.}
The residual generic witness survives over the abelian field
\(F_*\): the cubic splitting field has maximal abelian subfield
\(K\), so its cubic value can be varied while fixing every
rational character and \(F_*\).
At cyclotomic residue \(-1\) this gives the roots
\((-\zeta_3,1,-\zeta_3)\), whose ratios are never \(-1\).
The four elliptic tests of AC-G75, retaining the transformed
finite-torus summand in the fourth test, and the projective
Brauer argument AC-G78-type apply on the exact module of Step 2.
They give actual occurrence of the finite type
\(\tau_\vartheta=\operatorname{Ind}_{P_{2,1}}^{\mathrm{GL}_3(\mathbb F_q)}
(C_\vartheta^{(2)}\otimes1)\), with \(\vartheta\) regular of
order \(p\).  Positivity uses the residual genericity of every
constituent: spherical generic terms have equal Steinberg and
trivial multiplicities, nonspherical terms have nonnegative
difference, and the actual Steinberg term has positive difference.
It is not inferred from Steinberg presence alone.

The fixed full-\(q\)-level ordinary module is projective over
the weight algebra times \(\mathcal O[\mathrm{GL}_3(\mathbb F_q)]\).
Its \(\tau_\vartheta\)-multiplicity module is finite free over
the weight algebra and commutes with arbitrary weight base
change by the symmetric group-algebra identity of AC-G87.
A nonzero fiber therefore supplies the actual type at the
chosen classical weight \(e_*=133\), with the same finite
torus character and old projectors.  This makes the classical
weight assertion explicit.

The local parameter at \(q\) is a tame irreducible two-dimensional
block and an unramified line, with \(N=0\).
The block has Frobenius roots \(a,-a\); reduction against
\((-1,1,-1)\) forces the line to reduce to \(-1\).
Since \(\bar\gamma^2|_{I_p}\) has order \(134\), the middle
character is the unique individually conjugate-self-dual
residual global line.  A global \(2+1\) decomposition would
force its line to reduce to that middle character, contradicting
the local value.  Three lines cannot contain the irreducible
local block.  The genericity criterion excludes the remaining
noncuspidal discrete alternatives.  Thus the rank-three seed
is globally absolutely irreducible and its strong base change
is cuspidal.

At \(p\) the full-torus comparison has finite inertial tuple
\((\chi^j,1,\chi^j)\) and slopes \((-fe_*,0,fe_*)\),
\(f=538\).  The unequal middle character prevents a monodromy
link, and the equal endpoints have gap \(2e_*f\ne f\).
Hence \(N=0\).  The parallel weak-admissibility and descent
argument of AC-G79 gives the full ordered flag with quotients
\(u_1\chi^jc_p^{e_*},u_2,u_3\chi^jc_p^{-e_*}\), where
the \(u_i\) are unramified.  The old compact types give
actual away inertia \((\Psi,1,\Psi^{-1})\), not only its
semisimplification.

The ray-class exact sequence for \(E\), (G138.14), and
(G138.17) gives a totally real dihedral field \(D/\mathbb Q\),
with \(D/E\) cyclic of degree exactly \(p\), inertia \(p\)
at the prime over \(q\), and no other relative finite
ramification.  Indeed \(h_E=1\), \(p\Vert q+1\), and all
global units have trivial image in the local \(C_p\); the ray
\(p\)-part is exactly this \(C_p\), with conjugation inversion.
Odd degree precludes ramification at real places.

Here \((2/p)=-1\).  The Frobenius of \(p\) in the dihedral
field is a reflection, since its image on \(E\) is nontrivial.
It has order two, so \(p\) splits completely in \(D/E\).
Every \(D\)-completion over \(p\) is precisely the unramified
quadratic extension of \(\mathbb Q_p\), already contained in
the degree-\(538\) completion of \(F_*\).
Consequently adjoining \(D\) does not enlarge that completion.
This is the actual replacement for the older sufficient
assumption that the real quadratic field split at \(p\).

Set
\[
 H_0=F_*D=F_1TD,\qquad F_0=H_0^+,\qquad
 H_{\rm ab}\subset H_{\rm ab}T\subset H_0.
 \tag{G138.18}
\]
The intersection \(F_*\cap D=\mathbb Q\): every abelian
quotient of the dihedral group factors through \(E\), whereas
\(E\not\subset F_*\).  Thus \([H_0:\mathbb Q]=8p^2\),
and the two displayed indices are \(p\).
The extension \(H_0/F_0\) is a base change of
\(F_1/F_1^+\), finite-unramified and relatively split at
\(p,3,5\).  At \(q\) the unramified quadratic extension from
\(E\), followed by the tame inertia index \(p\), kills both
new finite characters.  At \(q_{\rm aux}\) they were already
killed by \(T\).  All remaining away inertia is exactly
\((\Psi,1,\Psi^{-1})\), with \(N=0\).

Rank-three irreducibility persists through the quadratic and
cyclic-\(p\) steps, since a self-twist would have cube one.
For rank two the added quadratic \(E\) cannot cause a split:
it ramifies at two, where \(V|_{G_{F_*}}\) is unramified.
Its following cyclic-\(p\) restriction again stays irreducible.
The retained cyclic base change and definite-unitary descent
therefore give both actual automorphic seeds on \(H_0\).
Their coefficient-prime flags restrict on the unchanged full
local field of degree \(538\).

\proofstep{5. The four new Kummer sectors and the strict rank gap.}
Here \(\mathcal S(L,\beta)\) denotes the finite residual
\(k(\beta)\)-valued strict group, zero in full local
cohomology at every \(p\)-place and on inertia at every
other finite place.  Write \(a_{L,0}\) for its dimension
and \(r_{\alpha,L}\) for the corresponding characterwise
cyclotomic \(k[[T]]\)-rank with the conventions of
AC-G98/AC-G100.
The four characters of \(H_{\rm ab}/K\) are still the restrictions
of \(1,\tau_{-11},\tau_8,\tau_{-88}\).
Set
\(\bar\alpha=\bar\chi^{-j}\bar\Psi^{-1}\bar c_p^{-133}\).
Using \(s\), not the later quartic field parameter, as the
quadratic-host index, the exact full-local Shapiro decomposition is
\[
 \begin{aligned}
 \mathcal S(H_{\rm ab},\bar\alpha)
   &=\bigoplus_{s=1,-11,2,-22}
                    \mathcal S(K,\bar\alpha\tau_{D_s}),\\
 &D_1=1,\quad D_{-11}=-11,\quad D_2=8,\quad D_{-22}=-88 .
 \end{aligned}
 \tag{G138.19}
\]
It includes the away conditions at every ramified quadratic place.
Now \(-133-1=-134\equiv134\pmod{268}\).
Since \(p\equiv1\pmod4\), the character
\(\tau_p=\bar c_p^{134}\) is that of \(\mathbb Q(\sqrt p)\).
Thus the character to split for Kummer descent, and its inverse
eigenspace character, are respectively
\[
 \begin{aligned}
 \eta_s&=\bar\alpha\tau_{D_s}/\bar c_p
       =\bar\chi^{-j}\bar\Psi^{-1}\tau_p\tau_{D_s},\\
 \eta_s^{-1}&=\bar\chi^j\bar\Psi\tau_p\tau_{D_s}.
 \end{aligned}
\]
For \(d=ps\) the exact splitting field is
\[
 \begin{aligned}
 M_d&=K(t,w_d),\qquad t^3=269,\qquad L_d=\mathbb Q(t,w_d),\\
 w_d^4+5dw_d^2+5d^2&=0,\\
 d&\in\{269,-2959,538,-5918\}.
 \end{aligned}
 \tag{G138.20}
\]
Indeed \(w_d=\sqrt d(\zeta_5-\zeta_5^{-1})\) cuts out the
quartic character \(\Psi\tau_p\tau_{D_s}\), and the faithful
cubic factor supplies \(K(t)\).  Its inverse has the same
splitting field.  These are exactly the fields of AC-G133 and
AC-G137; the additional \(\tau_p\) has not been discarded.
Their ramification at \(p\) belongs to this auxiliary coefficient
splitting cover, not to the actual seed's finite smooth \(p\)-type.

Their nontrivial cubic class sectors vanish by those exact
theorems.  To recall the reflection step,
\(\mathrm{ext}_{L_d}^{M_d}\circ\mathrm{Nm}_{M_d/L_d}=1+s_0\)
on the upstairs class group, with \(s_0\) a reflection.
Since \(\operatorname{Cl}(L_d)[p]=0\), that reflection acts
as minus one on the upstairs \(p\)-part.
Every conjugate reflection does likewise, and their order-three
products act trivially.  This kills precisely the nontrivial
cubic sector needed here.  It does not claim the whole upstairs
class group is prime to \(p\).

Full local triviality at every \(p\)-place, together with
away unramifiedness, puts Kummer representatives in the
all-finite-valuation-divisible group.  Its exact unit/class
sequence shows that the target strict group is a subspace
of the target unit eigenspace.  For
\(\Delta=\operatorname{Gal}(M_d/K)=C_3\times C_4\), the
complex places form one regular \(\Delta\)-set.  Equivariant
Dirichlet units have rational character
\(\mathbb Q[\Delta]-1\); the prime-to-\(p\) idempotents
preserve the same ranks modulo \(p\).
Each chosen nontrivial cubic and quartic eigenspace therefore
has dimension one after the finite value-field extension.
This also covers the real quartic fields
\(F_{-2959},F_{-5918}\); no CM assertion about their full normal
closure is needed.
There is no unit \(p\)-torsion: \([M_d:\mathbb Q]=24\),
whereas a primitive \(p\)th root of unity has degree \(268\).

For \(d=269\), AC-G134's actual full-local unit computation
gives strict zero for both cubic labels with this
\(\Psi^{-1}\) coefficient.  That computation uses both
ramified \(p\)-completions of absolute ramification degree six
and residue degree two; the first-residue degree-three
calculation is not substituted.
For each remaining \(d\), the preceding exact class and unit
argument gives at most one, with no unverified local unit
image asserted.  Hence
\[
 a_{{\rm ab},0}\le0+1+1+1=3<4=[H_{\rm ab}^+:\mathbb Q],
 \qquad r_{\alpha,{\rm ab}}\le3.
 \tag{G138.21}
\]
The local cubic inertia is nontrivial, so the permutation
correction in the cyclotomic control formula is zero.
Both cyclic steps in (G138.18) are unramified at \(p\),
and hence disjoint from the totally ramified cyclotomic
tower.  With the full prime-to-\(p\) residual splitting
cover, AC-G100's characterwise ascent applies at both steps:
\[
 r_{\alpha,H_0}\le3p^2=\tfrac34[F_0:\mathbb Q].
 \tag{G138.22}
\]
This is normalized rank transport, not a cubic-sector
\(\mu=0\) assertion.

\proofstep{6. Selected-place classes, filling, and one high host.}
Let \(\Omega\) be the original prime-to-\(p\) residual
splitting field containing \(K,\chi,\Psi,\bar c_p\).
It is ramified only at \(3,5,p\).
The maximal prime-to-\(p\) quotient of the initial normal
host has fixed field \(H_{\rm ab}\).  Each nontrivial
quadratic extension of \(H_{\rm ab}/K\) ramifies at two
or eleven, where \(\Omega/K\) is unramified.  Hence
\[
                         H_0\cap\Omega=K.
 \tag{G138.23}
\]
The normal closure of a later selected-place cyclic
\(p\)-extension, and every cyclotomic layer, has \(p\)-power
degree over this host, so this intersection persists.

Apply AC-G92 on \(F_0\) to a common prime-to-\(p\) cover
splitting the cubic tuple and the rational characters.
Choose at least one selected place.  The resulting
\(T'/F_0\) is cyclic of \(p\)-power degree and splits
every \(p\)-place.  The actual AC-G93 quotient supplies
a nonzero strict adjacent class, a nonzero central anti
class, and the rational sectors
\(\bar\gamma^{-1}\eta\), \(\bar\gamma^{-2}\eta\).
They are nontrivial since \(\bar\gamma(c)=1\) and
\(\eta(c)=-1\).  At \(p\), the cubic order three cannot
cancel a rational inertia order dividing \(268\); the
endpoint square has order \(134\).  All necessary local
\(H^0\) groups therefore vanish and stay zero after the
\(p\)-power restrictions.

The exact dual of the mixed obstruction is
\[
 \lambda=\eta\bar\Psi^2\bar c_p^{267}
         =\chi_{-15}\bar c_p^{267}.
 \tag{G138.24}
\]
Here \(2\cdot133+1=267\), so this is the same even dual
residual character as in the first residue.  It is even,
nontrivial on \(p\)-inertia, and has a totally real abelian
splitting field.  After adjoining that field
and \(\mu_p\), the real host is a finite \(p\)-extension
of an abelian field, including the selected extension.
The exact away-kernel/full-local duality of AC-G90--AC-G93,
with the retained abelian Iwasawa theorem and finite-\(p\)
ascent, gives bounded full dual \(H^1\).
Its injective restriction maps are eventually isomorphisms,
and corestriction is then zero modulo \(p\).
Perfect duality kills the mixed strict \(H^2\) obstruction
at a further layer.  The nonzero strict classes persist,
so the polarized filling and actual normalized Low paths
are supplied, rather than assumed.

The selected extension is disjoint from the cyclotomic tower,
since it splits \(p\).  AC-G100 preserves the ratio in
(G138.22).  The central and trivial sectors are rational;
their splitting fields are abelian before these finite
\(p\)-steps, so their dimensions are bounded by the same
retained rational-field argument.  On successive cyclotomic
layers, after a harmless reindexing,
\[
 a_n=r_\alpha p^n+\sum_i\min(t_i,p^n),\quad
 h_n=O(1),\quad z_n=O(1),\quad
 d_n=d_0p^n,\quad r_\alpha\le\tfrac34d_0.
 \tag{G138.25}
\]
Reindexing multiplies \(r_\alpha,d_0\) by the same power of
\(p\).  Thus eventually
\[
                     h_n+a_n+z_n+1<d_n,\qquad d_{v,n}>5.
 \tag{G138.26}
\]
Choose one layer beyond these and the filling thresholds;
the rational-frame thresholds below can be met at the same
time.  Denote it \(H/F\).
All subsequent steps are cyclic of \(p\)-power degree.
The two seeds stay irreducible and cuspidal by the determinant
self-twist argument and cyclic base change, with their exact
restricted flags and finite types.
No effective first layer or cubic \(\mu=0\) theorem has
entered this construction.

\proofstep{7. The actual rational frame and scalar auxiliary place.}
We recheck AC-G106/G121 for this nonquadratic character.
One has \(\bar\gamma(c)=\bar\Psi(c)(-1)^{133}=1\).
At \(p\), \(\Psi\) is unramified, so
\(\bar\gamma^2|_{I_v}=\bar c_p^{266}\) has order \(134\).
On the full inertia kernel of \(c_p\), \(\bar\gamma\)
is trivial; this is the literal \(\mathrm{Psi\_N}\) clause.
The characters \(\bar c_p\bar\gamma^{\pm2}\eta\) have
inertia exponents \(1\pm266\equiv267,3\pmod{268}\), both
nonzero.  These give the local separation and \(\mathrm{f2c}\)
clauses without any use of \(\Psi^2=1\).

The rational sectors of Step 6 give both positive classes
required by the source.  Their coefficient splitting fields,
and those of the three rational upper-bound sectors, are
abelian over the initial abelian host.  The abelian Iwasawa
theorem and finite-\(p\) ascent make their dimensions bounded.
The full arithmetic Euler and finiteness fields are AC-G96,
for arbitrary finite modules and local stable subspaces;
the local generators and class exponent are AC-G99.
Thus a constant \(C_0\) bounds the source rational strict
and trivial dimensions.  Choose the common high layer also
so that \(3C_0+7<d=[F:\mathbb Q]\) and every \(d_v>5\).
This verifies every rational \(\mathrm{Frame.Ok}\) and
\(\mathrm{Frame.Big}\) field with the actual full local maps.

There is an actual scalar auxiliary place.  The conductor-five
quartic field is disjoint from the joint cubic and conductor-\(p\)
fields by ramification at five.  The cubic splitting field has
maximal abelian subfield \(K\), disjoint from the conductor-\(p\)
cyclotomic field.  Hence the compatible order-two element
\[
                 \chi=1,\qquad\Psi=-1,\qquad\bar c_p=-1
 \tag{G138.27}
\]
exists in \(\operatorname{Gal}(\Omega/K)\) and has
\(\bar\gamma=1\).  By (G138.23) it extends fixing the host
normal closure, including the selected and cyclotomic steps.
Chebotarev gives a fresh rational prime completely split in
that closure, outside all bad places, and a residue-degree-one
place \(v_0\) with these values.  The source residual
characteristic polynomial is scalar there, while
\(\bar c_p(v_0)=-1\ne1\).  The determinant-only Sylvester
argument AC-G89 gives \(\mathrm{N4}\) for every native point;
it does not assume that every residual matrix is scalar.

Take the literal \(\mathrm{HidaSetting}\) with
\(\mathrm{St}\) the places over three and five.
The CM extension is finite-unramified and split at these
places, at all coefficient-prime places, and at \(v_0\).
At two all relevant representations are unramified, so two
is not required in \(\mathrm{St}\).  The source
\(\mathrm{p\_unr}\) concerns the fresh rational prime below
\(v_0\), not \(269\).  Retained \(\mathrm{LevEx}\) gives
the actual level with \(S'=S\).  All coefficient-prime maps
use full decomposition groups, not their global images.
Every global and local host degree has prime factors only
two and \(269\), so no such field contains \(\mu_{269}\),
whose degree has the factor \(67\).
Local reciprocity identifies the source \(\mathrm{mu}\),
the preimage of the \(p\)-primary torsion in the local
abelianization, with the closed commutator.  Thus
\(\mathrm{th}=1\) is valid.

\proofstep{8. The genuine rank-two point and its entire native component.}
Choose a finite host coefficient field containing the two
seeds, \(\mathcal O_0\), the finite characters, \(\kappa_2\),
and every embedding required by native \(\mathrm{Large}\).
It may be ramified; no condition on its ramification is
transferred to \(\mathcal O_0\).
For an actual split residual rank-two lattice, compactness
first gives a stable lattice.  The conjugation idempotents
split it into two rank-one eigenspaces, with
\(V(c)=\operatorname{diag}(-1,1)\).
Its residual invariant line is one of these eigenspaces,
so one off-diagonal family is divisible by the uniformizer.
After adjoining its square root, rescale one basis vector
to make both off-diagonal families divisible by the new
uniformizer.  The diagonal entries and conjugation matrix
are unchanged.  The result has actual split reduction
\(\eta\bar\gamma\oplus\bar\gamma^{-1}\).

On \(G_H\), \(\det V=1\), and direct matrix algebra gives
\[
 V((cgc^{-1})^{-1})=J_2V(g)^tJ_2,
 \qquad J_2=\begin{pmatrix}0&1\\1&0\end{pmatrix}.
 \tag{G138.28}
\]
Its difference with \(V(g)\) is diagonal, giving the exact
global normalization at the host witness.
The full local flags reduce to the labelled
\(\bar\gamma,\bar\gamma^{-1}\) lines.
Its actual away inertia \(\Psi\oplus\Psi^{-1}\), with
\(N=0\), kills the closed normal subgroup \(N_\gamma\)
generated by all \(I_w\cap\ker\bar\gamma\), \(w\nmid p\),
since reduction is injective on these order-four characters.
This subgroup is conjugation-stable.  The representation
therefore lies on \(\Gamma_\gamma=G_F/N_\gamma\), with
CM kernel \(G_\gamma=G_H/N_\gamma\).
This host assertion replaces the unavailable original
rational \(\mathrm{Kill}\); it is proved only after restriction.

The complete rational \(\mathrm{space\_ex}\) construction
gives \(B_2\subset A_2\subset R_2\), representing every
complete-local normalized ordinary test ring with these
full local maps.  The point above is such an actual test
point.  The native CM kernel maps onto \(G_\gamma\),
preserving conjugation and the full cyclotomic character.
The fixed \(\kappa_2\) of AC-G106 satisfies
\(\kappa_2\kappa_2^c=c_p^{-1}\), is unramified outside
\(3p\), and has finite away inertia of order dividing six.
It need only descend to the native group.  Take
\[
 \widetilde r_2=\kappa_2(r_2\circ\pi),\qquad
 \widetilde\lambda_i=\kappa_2\lambda_i,\qquad
 \bar t_i=\bar\kappa_2\bar\gamma^{1-2i}\quad(i=0,1).
 \tag{G138.29}
\]
These give the literal \(\mathrm{TwLink}\), \(\mathrm{Resid}\),
and \(\mathrm{Pair}\).  The determinant law is in the
completed invariant ring and is continuous by the
weight-zero atom bound of AC-G106-native.

Its joint away kernel can be larger than
\(I_w\cap\ker\bar\gamma\).  If both \(\bar t_i(g)=1\),
then \(\bar\gamma(g)^2=1\) and
\(\bar\kappa_2(g)=\bar\gamma(g)^{-1}\).
The universal matrix has finite prime-to-\(p\) order and
scalar residue \(\bar\gamma(g)\).  Factoring \(X^m-1\)
at that Teichmuller scalar makes its other factor a unit;
the matrix equals the scalar even over nonreduced rings.
The finite away \(\kappa_2\) value cancels it.  Thus the
literal \(\mathrm{Pair.I}\) holds, without assuming that
\(\Psi\) is quadratic.

The native map \(\Phi_2:\mathcal B_2\to B_2\) has closed
complete image containing traces and local characters.
For \(Q=r_2(g_0)-r_2(g_0^\star)=\operatorname{diag}(a,-a)\),
it contains \(a^2=\operatorname{tr}(Q^2)/2\).
The nonzero residual root and Hensel lifting put \(a\) in
the image.  Coordinate projectors and their word traces
give every diagonal entry and every
\(r_{01}(g)r_{10}(h)\), hence all of \(B_2\).
Thus \(\Phi_2\) is surjective on the actual completed ring.

Pull the fixed Hecke character \(Y\) of AC-G106 from \(K\)
to \(H\) and put
\(\Pi_2=\pi_V\otimes(Y|\cdot|^{1/2})\circ\det\).
The exact attachment normalization and
\(YY^c=|\cdot|^{-1}\) give the native CSD representation
\(V\kappa_2\).  Its weights are
\(b_\kappa(\tau)-133,b_\kappa(\tau)+133\), with all finite
phases retained.  They increase strictly, and no ratio
is \(c_p\): on open inertia the possible exponents would
be \(-1\) or \(\pm266-1\), all nonzero.
It is unramified at \(v_0\).  These actual arithmetic
\(\mathrm{PT}\) clauses make the retained \(\mathrm{ModU}\)
ideal \(J_2\) proper.

Native finite/torsionfree weight control makes
\(\mathcal B_2/J_2\) nonzero of component dimension
\(2d+1\).  The actual \(\mathrm{Setup}\), representation
point, \(\mathrm{Unr0}\), determinant-only \(\mathrm{N4}\),
and surjective \(\Phi_2\) are the inputs of
\(\mathrm{seedOf}\).  It recovers an entire component
with \(\beta\Phi_2=q\), hence genuine \(\mathrm{SeedIn}\).
Modulo the image of \(J_2\) and the weight ideal, every
inertia character equals its finite residue, even in the
nonreduced quotient.  The finite native weight module
therefore gives \(\mathrm{FinLam}\).
The actual \(\mathrm{PointEx}/\mathrm{NiceOf}\) inputs
use the verified rational \(\mathrm{Psi\_N}\) and give
\(\mathrm{P1'}\).  The complete rank-two component theorem
gives \(\mathrm{AllGood}\) for this proper image ideal.

The exact even-corner isomorphism AC-G106 keeps the
independent middle-character factor.  Its scalar lemma
gives inertia descent after removing \(\chi^j\), and
its local cancellation uses an inertia element with
\(\bar\gamma=1\) and \(\chi\) of order three, present on
all these hosts.  Its finite-support theorem supplies
one uniform \(M_E\subset B\), below every even
off-triangular prime, with
\(\sqrt{M_E+\lambda_\theta}=\mathfrak m_B\).
The middle character has not been set equal to one.

\proofstep{9. The new cubic native component and global support.}
Let \(N_\theta\) be the closed normal closure of all joint
residual away inertia kernels.  For (G138.4) it is generated
by \(I_w\cap\ker\chi\cap\ker\bar\gamma\): if both endpoints
are one, their ratio gives \(\bar\gamma^2=1\), and order three
forces both factors to be one.  Form the actual quotient
with full local maps.  A separating \(p\)-inertia witness
is obtained by cubing to kill \(\chi\) while retaining
order-\(268\) rational separation.
AC-G94 gives the complete universal functor, AC-G95 its
actual completed and algebraized rings, AC-G96 the
arbitrary-module Euler fields, and AC-G99 the generators
and class exponent.  Step 6 supplies actual lower points
and paths on this chart.

For the rank-three native Pair use \(\kappa_3=c_p^{-1}\),
not \(\kappa_2\); its conjugate norm is \(c_p^{-2}\).
The determinant law of \(\kappa_3r\), all three local
characters, and finite prime-to-\(p\) joint-kernel rigidity
give the actual native Pair and Resid.  The invariant
trace/projector argument AC-G102 gives the surjection
\(\Phi_3:\mathcal B_3\to B\).
The rank-three seed defines its arithmetic point with
weights \((1-133,1,1+133)\), the full finite phases, and
no cyclotomic ratio.  Its integral determinant law has
the required residual determinant; no claim of split
matrix reduction for an arbitrary rank-three stable
lattice is needed.  Its exact away types kill the
joint kernel and it is unramified at \(v_0\).
Thus \(\mathrm{PT}\) makes the native \(J_3\) proper.

The native quotient is finite torsionfree over the actual
weight algebra, of dimension \(3d+1\), with
\(\mathrm{Unr0}\) and the native prime-core properties.
AC-G107 recovers a whole component in the normalized
chart, with the literal identity \(\beta\Phi_3=q\).
It supplies a Type I prime \(\xi\), \(Z=\xi\cap B\), with
\[
                    I\subset Z,\qquad\dim B/Z=3d+1.
 \tag{G138.30}
\]
There is no inference that \(B/I\) is torsionfree.
The actual finite-weight image, native prime conditions,
and auxiliary \(\mathrm{N4}\) are AC-G102/AC-G104 on
these data.  The local presentations, normalized collision
bounds, and component chains of AC-G101/AC-G105 use
the separation and full Euler fields checked above.
Their dimension premise is (G138.26).  The lower core
is connected by the actual filled paths, and (G138.30)
is the strong seed.  Together with Step 8's uniform
even support, these give every premise of AC-G108.
It proves (G138.5), using the native image \(I\) for
Type I primes and \(I_{\rm off}=I\cap M_E\) for all
off-triangular primes.  The two ideals are not substituted
for one another.

These host rings and residual records depend on
\(\bar\gamma\), \(\mathrm{th}=1\), and the full local maps.
The residual tuple is constant for
\(e\equiv133\pmod{268}\).
The full unit-torus finite character is also constant:
the parallel algebraic weight change factors through
residue norms to \(\mathbb F_p^\times\), and \(e-133\)
is divisible by \(268\).
The ordinary functor leaves principal-unit powers
variable.  Thus the single seed at \(133\) supplies
support throughout the progression.  This does not
identify full characteristic-zero characters at different
weights or their original selected local ideals.

\proofstep{10. Original curves and original Type I points.}
Take an actual characteristic-\(p\) complete-DVR curve
from \(S_j\).  Embed its finite residue field and the
host residue field \(k\) compatibly in a finite \(k''\).
For \(V_0=k'[[z]]\), use
\(V_0\hookrightarrow V'=k''[[z]]\) and the host test
algebra \(C=k+zV'\), with coefficient map
\(\mathcal O\to k\).
The restricted matrices and their inverses lie in \(C\).
Only the representation is transferred; no scalar
extension or identification of the original deformation
ring is assumed.

The original exact modified Kill condition kills every
generator of \(N_\theta\) after restriction.
Its own selected-component AC-G65-flag gives the full
integral ordered inertia characters
\[
       \chi^j\Psi c_p^e,\qquad1,\qquad
                     \chi^j\Psi^{-1}c_p^{-e}.
 \tag{G138.31}
\]
The finite phase at \(p\) is \((\chi^j,1,\chi^j)\),
since \(\Psi\) is unramified there.
On this characteristic-\(p\) curve the fixed integral
characters reduce to constant residual values.
The paired labelled Hensel projectors descend the full
flags to \(C\); strict \(J\)-orthogonal normalization at
the new witness preserves the ordered identities and
characters.  The actual AC-G94 functor gives \(R\to C\),
annihilating \(\lambda_\theta\).

If the restricted curve were off both triangular loci,
(G138.5) would make its invariant image the closed
residue point.  Opposite-sign surviving entries \(x,y\)
would give a nonzero balanced monomial \(x^b y^a\) in
its maximal ideal, a contradiction in the domain \(C\).
The curve is triangular.  Its diagonal characters,
after removing their residues, are everywhere unramified
principal-unit characters of finite class-group exponent.
Principal units of a characteristic-\(p\) domain have no
torsion, so these characters are one.  Host characteristic
polynomials are the constant seed polynomials.
Every original \(g\in G_K\) has a positive power in
\(G_H\); its eigenvalues therefore satisfy finite-field
equations.  The characteristic coefficients are algebraic
constants, uniquely specified by their reductions.
This permits \(p\) to divide the host index.
The generic shape and trace-pairing argument AC-G109,
followed by the literal integral finite criterion over
\(\mathcal O_0\), gives original \(\mathrm{M2}\), including
vertical components.

For an original characteristic-zero Type I point, choose
a finite common value field and take the inverse image
of \(k\) in its integers as host test ring.
The same exact kernels, its own flags (G138.31), and
strict normalization give an actual point of \(A\).
The all-open and full-cyclotomic irreducibility proof
AC-G110 applies: on any open inertia subgroup killing
the finite characters, the different flag ratios are
the nonzero powers \(c_p^{\pm e},c_p^{\pm2e}\).
No quadratic identity for \(\Psi\) enters.
The point is Type I, so the first line of (G138.5)
makes its native point kill \(J_3\).
Its determinant and local characters are exactly those
of the restricted original representation times
\(c_p^{-1}\), with its own weights
\((1-e,1,1+e)\) and finite phases.
The native \(\mathrm{end2}\) direction gives actual
automorphy; the retained polarized adjoint theorem
gives host rigidity, which restriction/corestriction
descends to the original rational adjoint group.
The host degree is invertible in characteristic zero.
This is the consequence of the actual (G110.2), not
a reversal of the seed \(\mathrm{PT}\) implication.

\proofstep{11. The original model and normalized count.}
Every two-sided original point is Type I or Type II.
AC-G128, retaining the original modularity and
\(e_{\mathcal O_0}=1\) hypotheses, excludes every Type II
point for this quartic character, \(N=10\), and the
progression (G138.1): \(2e+1\equiv267\pmod{268}\), exactly
the weight class in that exclusion theorem.
It uses the original exact inertia, primitive level \(75\), its primitive
level-fifteen twist, and the actual bad roots.
It does not exclude the auxiliary level-\(75q_{\rm aux}^2\)
seed after host restriction.
Thus all original two-sided points are Type I and rigid.

AC-G122's normalized orbit/dual-point argument now gives
literal \(\mathrm{M4}\): after a common grading-torus
normalization the two derivations are coboundaries in
the line spanned by \(\operatorname{diag}(1,0,-1)\).
Entry determination makes them linearly dependent;
undoing the nonzero weight scalars gives the source
quantifier over every two-sided point.
AC-G66 gives \(\mathrm{M1},\mathrm{M3}\); Step 10 gave
\(\mathrm{M2}\).

At \(11\), \(K\) is inert, \(\psi_5(11)=1\), and the
degree-two Frobenius has \(\chi=1\).  This is a good
place for the original \(N=10\) model even though the
support host can ramify there.
Its original seed trace is
\(b=121^e+1+121^{-e}>3\).
The degree-zero original trace difference vanishes at
the origin and is nonzero at every two-sided point:
actual Type I traces have complex absolute value at
most three for every conjugate.  The grading torus
fixes this function.  This is literal homogeneous
\(\mathrm{M5}\), without a mixed Type II Selmer premise.

The host is reached from \(K\) by quadratic and cyclic
\(269\)-power steps.  The exact original-field descent,
\(\mathrm{WitA}\), and unlinking of AC-G111 therefore
apply.  The original fixed inertia, not the auxiliary
seed level, gives principal compact level
\(269N=2690\) by AC-G112.
Neither \(q_{\rm aux}\) nor \(q\), nor the host
ramification at two or eleven, enters that original
conductor bound.
The actual \(\mathrm{Af.TrData}\) carrier has bound
three, its exact integrality/conjugation properties,
and cardinality at most \(C_1Y^{A_1}\).

Apply the unchanged \(\mathrm{LemmaM}\) with
\(v_{\mathcal O_0}=e_{\mathcal O_0}v_{269}\), the
all-cover torus/isomorphism comparison, and the actual
\(\mathrm{LemmaM'}\) strict injection over the same
\(\mathcal O_0\).  This gives (G138.6)--(G138.7).
The classifier norm uses the labels \(j\in J_0\) with
controlled inert value one, retaining duplicate labels
when both occur.  It is one nonzero rational integer,
and its exact bound is
\[
 \log|\mathcal N_{e,J_0}|
 \le C_1\#J_0Y^{A_1}(2e\log121+1)
 \le5C_1\#J_0Y^{A_1+1}\log11.
 \tag{G138.32}
\]
Together with the valuation factor two this proves
(G138.8), including \(20C_1\) for both labels in its
last line.  No host ramification index rescales these
original normalized lengths.
\end{proof}

\paragraph{Scope and geometric dictionary.}
The theorem supplies actual support and original model
conditions on this one quartic second-residue progression.
The named native, published, and generic algebraic inputs
remain, and the finite arithmetic uses exactly AC-G137's
three new sectors.  No original Type II endpoint or mixed
Type II vanishing is assumed.
The independent AC-G134 second-residue strict vanishing supplies one
residual class-sector estimate; it does not imply model
rigidity or native support by itself.
No complementary quartic label, unchanged
quadratic ideal, effective cyclotomic depth, all-prime
estimate, or original family-average budget is asserted.

The previously audited SGA~1, Expose~V, sections~4 and~7
dictionary interprets the first cyclic cover as killing
its own order-\(269\) character local system, and the
later dihedral cover as killing the newly imposed tame
type.  Ordinary charts and native modular support are
then constructed on the actual host.  Original curves
and points enter through their separately checked
restriction and normalization maps.
This does not identify the original and host deformation
rings, construct a new motive, or assert a new
cohomological comparison theorem.  No Lean execution
is claimed.
